\section{Conclusion}
\label{sec:conclusion}

We have presented \textbf{SynBADD}, a bidirectional, anatomy-guided imputation framework for longitudinal 3D brain MRI. SynBADD operates natively in deformation space rather than in intensity space: a frozen 3D variational autoencoder encodes per-timepoint stationary velocity fields into a compact latent representation, a target-aware cross-attention transformer aggregates these latent codes across the longitudinal sequence, and a 3D U-Net refinement network conditioned on subject metadata via dynamic affine feature-map transforms (DAFT) produces the final SVF. Synthesis is completed by integrating the predicted SVF and warping the baseline scan with a differentiable spatial transformer, ensuring volumetrically coherent and topologically valid outputs by construction.

Across the ADNI-1 cohort, SynBADD substantially outperforms two recent baselines---NODER and ReMiND---and a simple Forward Difference imputation. Improvements are consistent across all five cross-validation folds and across the eight Alzheimer's-disease-relevant regions of interest evaluated, with PSNR and Dice gains particularly pronounced in the medial temporal lobe and ventricular system. The predicted deformation fields are diffeomorphic and exhibit a mean Jacobian determinant slightly below unity, reflecting the mild net tissue contraction expected in a cohort that includes subjects with MCI and AD. Subgroup analysis confirms that performance is broadly stable across diagnostic groups and target follow-up timepoints, indicating that the model accommodates the larger and more heterogeneous deformations associated with advanced neurodegeneration without requiring per-group specialisation.

Two controlled robustness experiments (Section~\ref{sec:robustness}) examined the model's behaviour under reduced and asymmetric input availability. The full UNet$+$DAFT model exhibits remarkable invariance to both the number and identity of available input timepoints---predictions are virtually unchanged whether all or only the baseline scan are provided. A controlled No-UNet$+$No-DAFT ablation, however, recovers the expected temporally-graded behaviour: imputation quality degrades smoothly and monotonically as inputs are dropped, with a weak but statistically significant recency effect indicating that proximal inputs contribute more than distant ones. Crucially, neither configuration shows a past-versus-future asymmetry. Because the ablation removes both U-Net refinement and DAFT, this contrast indicates that the combined refinement-and-conditioning pathway is a major contributor to imputation quality in the conditioned configuration. For incomplete longitudinal records, this is a useful property: the model degrades gracefully under sparse longitudinal records. Architecturally, it motivates a factorial ablation that separates U-Net refinement from DAFT conditioning before drawing conclusions about whether the temporal transformer is necessary, and an avenue for future work is to develop training regimes---for example, random input dropout that prevents the conditioning pathway from dominating---that engage both pathways more symmetrically.

\paragraph{Limitations.} Several limitations bound the present work. First, our evaluation is restricted to structural T1-weighted MRI from a single cohort (ADNI-1). Multi-modal extensions---incorporating, for example, FDG-PET or amyloid imaging---and validation on independent cohorts (ADNI-2, ADNI-3, AIBL) would be necessary to assess the generality of the conclusions. Second, the discrete timepoint structure of the ADNI-1 schedule means SynBADD operates on a coarse temporal grid; extending the model to handle continuous-time prediction without sacrificing the discrete-token attention formulation is non-trivial. Third, the SVF-based deformation framework assumes that anatomical change between consecutive timepoints is a smooth, topology-preserving warp. This assumption is valid for the slow atrophy patterns of AD, but breaks down in the presence of acute interventions, tumour growth, or surgical resections. Fourth, our subject metadata conditioning is limited to age and diagnostic label; richer biomarkers, such as APOE-$\varepsilon 4$ status or CSF tau levels, may further improve subject-specific adaptation. Finally, the temporal span of ADNI-1 (sc to m36) is short relative to the multi-decade course of Alzheimer's disease; longer-horizon projection remains an open problem that almost certainly requires per-subject regularisation beyond what a population-trained transformer can provide.

\paragraph{Future work.} We see three concrete directions. First, training regimes that explicitly engage the temporal pathway---random input dropout during training, or a contrastive objective that penalises the model for producing identical outputs under different input configurations---should recover the temporally-aware behaviour observed in the No-UNet$+$No-DAFT ablation while retaining the high absolute quality of the conditioned model. Second, downstream tasks such as longitudinal atrophy-rate estimation, conversion-to-AD prediction, and biomarker-trajectory analysis offer clinically meaningful evaluations that go beyond the per-scan image-fidelity metrics reported here. Third, integrating SynBADD into a probabilistic framework---for instance by treating the predicted SVF as a posterior mean and sampling from a learned uncertainty estimate---would enable uncertainty quantification on the imputed scans, an important property to evaluate before considering clinical decision support.

In summary, SynBADD demonstrates that operating natively in deformation space, combined with bidirectional temporal attention and clinical-metadata conditioning, yields a longitudinal imputation framework that is both quantitatively strong and anatomically interpretable. The honest examination of its component contributions presented here---through controlled ablations and statistical-significance testing---provides a transparent foundation upon which future work can build.
